\section{Introduction}
\label{sec:intro}

Reinforcement learning from an LLM judge has a default that is rarely stated: the judge is shown
the candidate turn, in the context that preceded it, and scores that. For single-turn tasks this
is the whole story; for multi-turn goal-oriented dialogue it assumes that a turn's value is visible
at the moment it is produced.

Motivational interviewing (MI) is a setting where the assumption is plainly false. MI is a
counselling style for helping people resolve ambivalence about change, and almost every technique
it teaches is defined by its downstream effect: a reflection is good if the patient elaborates, an
open question is good if it surfaces change talk, and premature reassurance is bad because of what
it forecloses two turns later \citep{miller2013mi}. A judge shown only the turn sees none of this;
we hypothesise that it rewards what is locally visible, above all whether the patient felt
validated, which a policy can exploit with unearned praise.

This paper presents \textbf{\LAGRPO}. Before the judge scores a candidate therapist turn, the
patient simulator and the current policy roll the conversation forward $K$ further turns, and the
judge scores the extended transcript (a turn is one utterance, so $K{=}5$ appends patient,
therapist, patient, therapist, patient). The lever itself is not new:
\citet{baruch2025pto} introduced $K$-turn look-ahead for this task inside preference-tree
optimisation (PTO), where the look-ahead scores select a chosen and a rejected turn for DPO. We
move it to Group Relative Policy Optimization \citep[GRPO;][]{shao2024deepseekmath}, where a
single rollout is a one-sample estimate of where a turn leads and every candidate, not a filtered
pair, is trained on (\S\ref{sec:method}). We ask whether the look-ahead reward retains its
benefit under group-relative optimisation and, in more depth, what it does to the therapist the
policy becomes.

We answer with a controlled pair: two ten-iteration GRPO runs, identical in therapist,
simulated patients, oracle, hyperparameters and personas, whose tracked configurations differ
only in whether the oracle scores the candidate turn alone ($K{=}0$) or with its five-turn
continuation ($K{=}5$). Every model state is scored on eight MI instruments by the training oracle
and by a held-out judge from a different model family, and every utterance of every evaluation
conversation is coded with an MI process coder, so that the analysis can follow what each
therapist turn is and what the patient says next.

The reward result is large and consistent: look-ahead leads on the rewarded rubric from the
fourth iteration on under both graders, leads on all eight instruments at the matched endpoint,
beats the turn-level policy's \emph{best} checkpoint as well as its last, and survives an
independent re-draw of the evaluation set (\S\ref{sec:reward}). The process analysis is the
paper's centre. Under turn-level reward the policy converges on non-specific praise, the
MI-inconsistent behaviour of affirming the patient regardless of what they said, delivered most
often after the patient's \emph{sustain} talk; a judge-free lexical marker corroborates the coded
rate (\S\ref{sec:therapist}). Under look-ahead it converges on complex reflections, delivered
after the patient's \emph{change} talk; a look-ahead reflection is followed by change talk in
$85$--$89\%$ of cases, and patient change talk keeps rising through the session where under
turn-level reward it falls back late (\S\ref{sec:patient}). Both graders agree on that picture
and disagree on how much the look-ahead policy still praises.

Our contributions are: (i) \LAGRPO, the look-ahead reward of \citet{baruch2025pto} moved from
preference-tree DPO to a group-relative optimiser, evaluated in a controlled two-grader design
with a replicated endpoint; and (ii) an utterance-level account of what each reward horizon
teaches, in MI's own terms: which therapist behaviours each policy learns and unlearns, what each
behaviour yields in the patient's next turn, and how change talk develops over a session.
Appendix~\ref{app:saturation} documents the training oracle's saturation at the winning
checkpoint, the reason the held-out judge carries the per-conversation claims.
